\documentclass[10pt,a4paper]{amsart}
\usepackage[margin=19mm]{geometry}
\usepackage{amsmath,amssymb,mathtools}
\usepackage[scr=rsfs,cal=euler]{mathalfa}
\usepackage{xfrac}
\usepackage[unicode,hidelinks,bookmarks=false]{hyperref}
\newcommand{\ie}{i.e.\ }
\newcommand{\cf}{cf.\ }
\newcommand{\resp}{respectively\ }
\newcommand{\Subsection}{\subsection}
\providecommand{\nobreakdash}{\nobreak\hbox{-}}
\setlength{\emergencystretch}{2em}
\multlinegap0pt
\usepackage{csquotes}
\usepackage{mathtools}
\usepackage{booktabs}
\usepackage{multirow}
\usepackage{siunitx}
\usepackage{float}
\usepackage{bm}
\usepackage{amssymb}
\newtheorem{theorem}{Theorem}
\title{Virtual Constraint for a Quadrotor UAV Enforcing a Body-Axis Pointing Direction}
\author{Alexandre Anahory Simoes, Leonardo Colombo, Juan Giribet, Efstratios Stratoglou}
\date{Source: arXiv:2602.24268v1 (CC BY 4.0)}
\begin{document}
\maketitle

\noindent\emph{Source and licence.} Virtual Constraint for a Quadrotor UAV Enforcing a Body-Axis Pointing Direction. Version arXiv:2602.24268v1 (CC BY 4.0), distributed by arXiv under the Creative Commons Attribution 4.0 International licence, http://creativecommons.org/licenses/by/4.0/, as recorded in the arXiv licence metadata for this exact version and shown on the abstract page https://arxiv.org/abs/2602.24268v1. Full licence notice: \texttt{licenses/arxiv-2602.24268-CC-BY-4.0.txt}.\par\smallskip

\noindent\textit{Editorial scope.} Counted unit: the article's theorem eq:tau23-laws (source0.tex lines 603--615). The article's complete original proof of it is delivered with it (source0.tex lines 617--692, one \texttt{proof} environment of 76 lines). No mathematics is written, repaired or re-derived: the statement and the proof are the article's own lines, in the article's own notation, and no retained line is substituted or re-flowed.  Retained uncounted as the source-local context the proof needs: lines 195--201; lines 535--537; lines 568--570.  Nothing was omitted from any retained span, so the package carries no omission record.  No retained line contains a citation, so the wrapper carries no bibliography and no bibliography entry.  No retained line contains a figure environment, an image-inclusion command or drawing code, so no image is delivered and the package declares no asset dependency.  Editorial formatting only, in addition: an AMS \texttt{amsart} preamble that restates the article's own declarations and macros used by the retained lines, namely the environments \texttt{theorem} on separate counters, together with the standard packages those lines need.  The retained environments print in this document as Theorem 1 in order of appearance, whereas the article prints the same items with the numbering of its own full document.  The complete native source of the article as distributed in the arXiv e-print for this exact version is bundled unchanged as \texttt{sources/195454/source0.tex}.
\begin{align}
\dot{p} &= v, \label{eq:transl_kin}\\
m \dot{v} &= f R e_3 - m g e_3, \label{eq:transl_dyn}\\
\dot{R} &= R \hat{\Omega}, \label{eq:rot_kin}\\
J \dot{\Omega} + \Omega \times J \Omega &= \tau,
\label{eq:rot_dyn}
\end{align}

\begin{equation}\label{eq:M1-def}
\mathcal{M}_1 := \{(R,p)\in \mathcal{M}:\ b_1 = -\hat{p},\ \ e_3^\top p=z_0\}.
\end{equation}

\begin{equation}\label{eq:Omega23-constraints}
\Omega_3 = \frac{v_2}{\rho},\qquad \Omega_2 = -\frac{v_3}{\rho}.
\end{equation}

\begin{theorem}
    Suppose that the control system \eqref{eq:transl_kin}-\eqref{eq:rot_dyn} is actuated by $(f,\tau_2,\tau_3)$. Then, the thrust controller
    $$f=\frac{mg}{e_3^\top R e_3},$$
    and the torque controllers given by
    
\begin{align}
\tau_2 &= J_2\Big(\Omega_1\Omega_3 + \frac{1}{\rho}\Big(-\frac{f}{m}+g s_3\Big)\Big)+(J_1-J_3)\Omega_1\Omega_3,\nonumber\\
\tau_3 &= J_3\Big(-\Omega_1\Omega_2 - \frac{g}{\rho}s_2\Big)+(J_2-J_1)\Omega_1\Omega_2,\label{eq:tau23-laws}
\end{align} are the unique controllers ensuring that the submanifold \eqref{eq:M1-def} is a virtual constraint for this control system, where
    \[
        s_2=e_3^\top b_2=e_3^\top R e_2,\qquad s_3=e_3^\top b_3=e_3^\top R e_3.
    \]
\end{theorem}

\begin{proof}
From $e_3^\top p=z_0$ we require $e_3^\top v=0$ and
\[
0=\frac{d}{dt}(e_3^\top v)=e_3^\top \dot v.
\]
Using the translational dynamics $m\dot v=f b_3-mg e_3$ yields
\[
0=e_3^\top\dot v=\frac{f}{m}e_3^\top b_3-g,
\]
hence, provided $e_3^\top b_3\neq 0$,
\begin{equation}\label{eq:f-law}
{\quad f=\frac{mg}{e_3^\top b_3}=\frac{mg}{e_3^\top R e_3}.\quad}
\end{equation}

To keep $b_1=-\hat{p}$ invariant, it suffices to enforce that the velocity-level constraints \eqref{eq:Omega23-constraints} are preserved by the dynamics. Differentiate \eqref{eq:Omega23-constraints}:
\begin{equation}\label{eq:diff-Omega23}
\dot\Omega_3 = \frac{d}{dt}\Big(\frac{v_2}{\rho}\Big),\qquad
\dot\Omega_2 = -\frac{d}{dt}\Big(\frac{v_3}{\rho}\Big).
\end{equation}
We compute the required right-hand sides.

First, note that $\dot\rho = \hat{p}^\top \dot p = \hat{p}^\top v = -b_1^\top v = -v_1$, hence
\begin{equation}\label{eq:rho-dot}
\dot\rho = -v_1.
\end{equation}
Next, since $v_i=b_i^\top v$,
\[
\dot v_i = \dot b_i^\top v + b_i^\top \dot v,\qquad \dot b_i=(R \Omega)\times b_1.
\]

Using $\dot v=\frac{f}{m}b_3-ge_3$ we obtain:
\begin{align*}
\dot v_2 &= ((R \Omega)\times b_2)^\top v + b_2^\top \Big(\frac{f}{m}b_3-ge_3\Big)\\
&= ((R \Omega)\times b_2)^\top v - g\, s_2,\\
\dot v_3 &= ((R \Omega)\times b_3)^\top v + b_3^\top \Big(\frac{f}{m}b_3-ge_3\Big)\\
&= ((R \Omega)\times b_3)^\top v + \frac{f}{m}- g\, s_3.
\end{align*}
A direct expansion in the body basis gives
\[
((R \Omega)\times b_2)^\top v =- v_1\Omega_3 + v_3\Omega_1,\quad\]
\[
((R \Omega)\times b_3)^\top v = v_1\Omega_2 - v_2\Omega_1.
\]
Substituting the on-manifold relations $v_2=\rho\Omega_3$ and $v_3=-\rho\Omega_2$ from \eqref{eq:Omega23-constraints} yields
\begin{align}
\dot v_2 &= -v_1\Omega_3 - \rho\Omega_2\Omega_1 - g s_2,\label{eq:vdot2-simpl}\\
\dot v_3 &= v_1\Omega_2 - \rho\Omega_3\Omega_1 + \frac{f}{m}- g s_3.\label{eq:vdot3-simpl}
\end{align}
Combining \eqref{eq:diff-Omega23}--\eqref{eq:rho-dot} with \eqref{eq:vdot2-simpl}--\eqref{eq:vdot3-simpl} gives the required accelerations:
\begin{equation}\label{eq:Omega23dot-required}
{\begin{aligned}
\dot\Omega_3 &= -\Omega_1\Omega_2 - \frac{g}{\rho}s_2,\\[1mm]
\dot\Omega_2 &= \Omega_1\Omega_3 + \frac{1}{\rho}\Big(-\frac{f}{m}+g s_3\Big).
\end{aligned}}
\end{equation}

Assume $J=\mathrm{diag}(J_1,J_2,J_3)$. Then
\[
J\dot\Omega+\Omega\times J\Omega=\tau,
\]
and the second and third components satisfy
\begin{align*}
\dot\Omega_2&=\frac{1}{J_2}\Big(\tau_2-(J_1-J_3)\Omega_1\Omega_3\Big),\\
\dot\Omega_3&=\frac{1}{J_3}\Big(\tau_3-(J_2-J_1)\Omega_1\Omega_2\Big).
\end{align*}
Imposing \eqref{eq:Omega23dot-required} gives the feedback laws
\begin{equation}
{\begin{aligned}
\tau_2 &= J_2\Big(\Omega_1\Omega_3 + \frac{1}{\rho}\Big(-\frac{f}{m}+g s_3\Big)\Big)+(J_1-J_3)\Omega_1\Omega_3,\\
\tau_3 &= J_3\Big(-\Omega_1\Omega_2 - \frac{g}{\rho}s_2\Big)+(J_2-J_1)\Omega_1\Omega_2,
\end{aligned}}
\end{equation}
together with $\tau_1=0$.

Finally, substituting the thrust law \eqref{eq:f-law} into \eqref{eq:tau23-laws} yields a closed form in terms of $(R,p,v,\Omega)$ only. 
\end{proof}

\end{document}
